% Lösungen Einheit 4 von 5 – Übergangsmodell (zusammenbau v0.1)
\einheitenkopf{Einheit 4 von 5 · Übergangsmodell}

\erg{15}{a) von A nach C \quad b) Schleife an A mit $0{,}9$, Pfeil A → B mit $0{,}1$, Pfeil B → A mit $0{,}15$, Schleife an B mit $0{,}85$ \quad c) $M = ((0{,}8 | 0{,}35), (0{,}2 | 0{,}65))$ \quad d) $0{,}4$: 40\,\% der Gäste von Q essen am nächsten Tag in P \quad e) $v_1 = (550 | 450)$, $v_2 = (575 | 425)$ \quad f) $(400 | 600)$ aus $0{,}7x + 0{,}1y = 340$ und $x + y = 1000$ \quad g) $a = 500$ (aus $0{,}8a + 60 = 460$) \quad h) Der Vorgänger wäre $x = -50$, $y = 1050$ (aus $0{,}9x + 0{,}2y = 165$ und $x + y = 1000$) – negative Kundenzahl, also unmöglich; dass die Summe 1000 stimmt, reicht nicht}
\erg{16}{a) Fehler: Zeilen und Spalten vertauscht – die Pfeile aus A gehören in die erste Spalte. Richtig: $M = ((0{,}6 | 0{,}1), (0{,}4 | 0{,}9))$ \quad b) Die Komponente j von $v_n$ wird mit der Spalte j multipliziert; diese Spalte verteilt den Bestand von Zustand j auf alle Zielzustände – also steht in ihr, wohin Zustand j abgibt \quad c) ja: $B_1 = 760$, $B_2 = 732$ – beide Monate über 700}
